\originalchapter{闭合恒星模型与三个时间尺度}
本章对应8.901第7--8次课，并将第\ref{ch:stellar}章的基本结构式扩展为可解边界问题。知道四条微分方程还不足以得到一颗恒星；需要物质关系、边界条件，以及解释何时可以忽略时间变化。

\originalsection{质量坐标下的四个结构方程}
以包围质量$m$为自变量，链式法则把式\eqref{eq:stellar}改成
\[
 \frac{\dd r}{\dd m}=\frac1{4\pi r^2\rho},\quad
 \frac{\dd P}{\dd m}=-\frac{Gm}{4\pi r^4},\quad
 \frac{\dd L}{\dd m}=\epsilon_{\rm nuc}-\epsilon_\nu-T\frac{\dd s}{\dd t},
\]
\[
 \frac{\dd T}{\dd m}=\frac{T}{P}\nabla\frac{\dd P}{\dd m}.
\]
这里$\dd s/\dd t$跟随同一质量元，$s$为单位质量熵。它保留收缩、储热或释放热的贡献；热稳态主序中才近似忽略。组成变化的核能与化学项需一致放在$\epsilon_{\rm nuc}$定义中，不能一方面用净核产热、一方面再重复计入同一结合能变化。

中心$r,L\to0$且$P,T$有限；表面需要总质量$M$、大气压强温度关系及光度半径的联系。求解通常调中心$P_c,T_c$使表面条件成立，或作两端射击／松弛。任意指定$M,R,L,T_c$一般过度约束。状态方程$\rho(P,T,X_i)$、不透明度、核反应率和组成演化负责闭合，不是由重力方程自动产生。

\originalsection{对流判据是在同一压强下比较密度}
小气团上移，声速使其很快和环境达到同压强，若移动比热交换快则近绝热。均匀组成、理想气体下上移$\dd\ln P<0$，气团$\dd\ln T=\nabla_{\rm ad}\dd\ln P$，环境$\nabla_{\rm env}\dd\ln P$。若环境梯度更大，环境更冷，气团较轻继续上升，故不稳定$\nabla_{\rm env}>\nabla_{\rm ad}$。

一般状态方程定义$\delta=-(\partial\ln\rho/\partial\ln T)_{P,\mu}$、$\varphi=(\partial\ln\rho/\partial\ln\mu)_{P,T}$、$\nabla_\mu=\dd\ln\mu/\dd\ln P$。含组成梯度时比较密度给Ledoux判据
\[
 \nabla_{\rm env}>\nabla_{\rm ad}+\frac{\varphi}{\delta}\nabla_\mu .
\]
向内增加的$\mu$可稳定分层；均匀组成恢复Schwarzschild判据。边缘稳定判据只告诉是否增长，尚未给实际热通量和混合速率，后者要对流模型。辐射计算所需$\nabla_{\rm rad}$超过阈值时说明纯辐射模型不自洽，不能继续照原温度梯度而声称已包含对流。

\originalsection{动力学、热与核时间区分了不同近似}
动力学时间$t_{\rm dyn}\sim\sqrt{R^3/(GM)}$来自重力加速度使位置改变$R$的时间。声穿越时间在静水平衡温度下同量级。若演化慢于此时标，星体可以每一刻近静水平衡。

热调整的Kelvin--Helmholtz时间$t_{\rm KH}\sim GM^2/(RL)$来自可释放引力与热能除以光度；均匀、理想气体维里模型的可用总能为$|W|/2$，精确系数取决于结构。核时间$t_{\rm nuc}\sim0.007fXM c^2/L$则由参与烧氢的质量分数$f$与氢丰度$X$决定。核燃料并非全星同时可用，所以不能仅以$0.007Mc^2$推所有恒星寿命。

普通长期烧氢星常有$t_{\rm dyn}\ll t_{\rm KH}\ll t_{\rm nuc}$，允许准静态的热平衡主序；核闪或坍缩时某个不等式失效，就需时间依赖的能量或动力学方程。脉动的周期靠近动力学时间，不意味着平均核演化也同样迅速。

\originalsection{同源尺度律解释光度质量关系的方向}
粗略同形模型有$\rho\sim M/R^3$、$P\sim GM^2/R^4$。气压主导给$T\sim\mu m_u GM/(k_BR)$。扩散光度量级$L\sim RT^4/(\kappa\rho)$（省略共同辐射常数），若$\kappa$近常数，
\[
 L\propto\frac{\mu^4 M^3}{\kappa}.
\]
光度随质量增长很快，因此大质量星虽燃料多，寿命仍短。Kramers$\kappa\propto\rho T^{-7/2}$时改为$L\propto\mu^{15/2}M^{11/2}R^{-1/2}$。两者是不同温区的近似，不是同时对全主序成立的两个精确经验公式。

若核产热每单位质量$\epsilon\propto\rho^uT^v$，另有$L\sim\epsilon M\propto\mu^vM^{1+u+v}R^{-3u-v}$，两套估计结合才可约束$R(M)$。$v$局部依温度变化，对流、大气和辐射压也会打破同源假设，所以这些尺度律不代替实际恒星演化计算。多方球则用$P=K\rho^{1+1/n}$闭合一族数学结构，PS4.3的实际曲线和解析$n=1$检验说明它的可检验性及模型边界。
